\documentclass[11pt]{article}
\usepackage[a4paper,margin=2.45cm]{geometry}
\usepackage{amsmath,amssymb}
\usepackage{booktabs}
\usepackage{graphicx}
\usepackage{microtype}
\usepackage{natbib}
\usepackage[hidelinks]{hyperref}
\usepackage{xcolor}
\usepackage{enumitem}
\usepackage{siunitx}
\usepackage{threeparttable}
\usepackage{caption}
\usepackage{subcaption}

\sisetup{round-mode=places,round-precision=4,detect-all}
\setlist{nosep,leftmargin=*}
\newcommand{\code}[1]{\texttt{#1}}
\newcommand{\dataset}{RELISH}

\IfFileExists{generated/results_macros.tex}{
  \input{generated/results_macros.tex}
}{
  \newcommand{\TrainQueries}{611}
  \newcommand{\CalibrationQueries}{339}
  \newcommand{\TestQueries}{2240}
  \newcommand{\CandidatePairs}{191245}
  \newcommand{\UniqueDocuments}{162971}
  \newcommand{\PrimaryGain}{pending}
  \newcommand{\PrimaryCI}{pending}
  \newcommand{\PrimaryPValue}{pending}
  \newcommand{\PrimaryCoverage}{pending}
}

\title{When should bibliographic metadata alter related-article recommendations?\\
A leakage-aware selective evaluation}
\author{Zichen Feng\\
Department of Computer Science and Technology, University of Sanya, Sanya, Hainan, China\\
\texttt{3353854381@qq.com}}
\date{}

\begin{document}
\maketitle

\begin{abstract}
Citation counts and publication years can alter scientific-paper recommendations,
but their apparent value is easily overstated by target leakage, post-publication
metadata, repeated use of test collections, and comparisons against weak content
retrievers.  We re-evaluate bibliographic intervention as a selective decision: for
each seed paper, either retain a fixed SPECTER2 ranking or choose one of eight
pre-specified citation/recency rerankings.  Before materialising locked-test labels, we
publicly froze the split, action family, metric definitions, calibration rule, and stopping
criterion.  The confirmatory experiment uses \dataset, an expert-curated biomedical
related-article benchmark with \CandidatePairs{} graded query--candidate judgements.
The policy is trained on \TrainQueries{} seed papers, calibrated on \CalibrationQueries{},
and evaluated once on \TestQueries{} locked seeds.  Comparators include BM25,
BGE-small, SciNCL, SPECTER2, a training-only LambdaRank fusion, a global metadata
action, local k-nearest-neighbour action selection, and a gradient-boosted gain
predictor.  We report mean NDCG@10 change, matched-coverage conditional harm,
negative shortfall, and a genuine selective risk--coverage curve.  At the locked
operating point, the BiblioGuard policy changes NDCG@10 by \PrimaryGain{}
(95\% paired-bootstrap interval \PrimaryCI{}; randomisation \(p=\PrimaryPValue\))
at \PrimaryCoverage{} coverage.  The accompanying audit shows why the earlier
SCIDOCS citation result is retrospective rather than confirmatory: current citation
counts are coupled to a citation-prediction target and the official test collection
had already influenced method development.  The contribution is therefore a
leakage-aware empirical assessment and reproducible decision protocol, not a formal
safety guarantee.
\end{abstract}

\noindent\textbf{Keywords:} scientific paper recommendation; bibliographic metadata;
selective prediction; risk--coverage; test leakage; reproducibility

\section{Introduction}

Related-article recommenders help researchers navigate a literature that is too large
to inspect manually.  Text encoders estimate topical similarity, while bibliographic
signals such as citation count and publication year encode visibility and recency.
Hybridising these signals appears attractive, but a citation prior can reward age,
field size, and accumulated attention instead of query-specific relevance.  A fixed
metadata weight can therefore improve a mean metric while degrading many individual
queries.

This paper asks a narrower question: \emph{given a strong content ranking, when does
the available development evidence justify changing it with bibliographic metadata?}
The question is a selective-decision problem because abstention has a concrete
meaning---retain the content ranking.  It is also an evaluation-design problem.
Citation metadata gathered years after a benchmark was created can leak future
information; on citation-prediction collections it can be directly coupled to the
target.  Cross-validation does not restore confirmatory status when the same official
test collection has already been used for feature, neighbourhood, or threshold
choices.

An earlier version of this work did not resolve those threats.  Its only material gain
occurred on SCIDOCS, where post-hoc citation counts were used for a citation-prediction
task.  It also varied several configurations on the same official tests, lacked
matched-coverage selectors and modern scientific encoders, and called an aggregate
gain/coverage plot ``risk--coverage'' without computing conditional selective risk.
We do not conceal those facts: the old four-dataset results are retained only as a
retrospective audit.

The revised study makes four contributions:

\begin{enumerate}
  \item A publicly timestamped protocol and physically separated
  \code{fit}, \code{calibrate}, \code{freeze}, and \code{evaluate} phases.  Frozen
  decisions are hashed before locked labels can be materialised.
  \item A new confirmatory evaluation on human graded relevance rather than a
  citation-prediction target, with exact-CorpusId metadata provenance and explicit
  missingness.
  \item Comparisons with recent scientific document representations and a supervised
  learning-to-rank fusion, plus selectors evaluated at identical coverage budgets.
  \item Query-level inference and risk reporting: paired effects, harm probability,
  severe harm, negative shortfall, selective risk--coverage, and AURC.
\end{enumerate}

Our stopping rule is deliberately adverse to a positive narrative.  If a fixed action,
LambdaRank, or a simpler gain predictor matches or exceeds BiblioGuard, that result is
reported and the method is not described as safer or more effective.

\section{Related work}

\subsection{Scientific-paper recommendation and representation}

Research-paper recommenders combine content, collaborative, and graph signals, but
their datasets and evaluation designs vary widely \citep{zhang2023survey}.  RELISH
provides more than 180,000 document-pair judgements from over 1,500 scientists
\citep{brown2019relish}; a broad comparison found that representation choice and
domain adaptation materially affect biomedical recommendation
\citep{zhang2022comparative}.  The current study uses RELISH's three graded relevance
levels and its fixed 60-document candidate pools.

Dense scientific representations have advanced beyond compact general-purpose
sentence encoders.  SPECTER uses citation links to learn paper representations
\citep{cohan2020specter}; SciNCL samples neighbourhoods in a citation-embedding space
\citep{ostendorff2022scincl}.  SciRepEval showed that scientific representations do not
generalise uniformly across task formats and introduced task-format adapters in
SPECTER2 \citep{singh2023scirepeval}.  We therefore compare lexical BM25, BGE-small,
SciNCL, and the SPECTER2 proximity adapter on exactly the same candidate pools.

Recent academic recommendation work uses large citation graphs, content, and learned
multi-stage systems \citep{stergiopoulos2024academic,kanwal2024multiple}.  Knowledge-
aware citation recommendation and OAG-Bench further illustrate the stronger graph and
benchmark comparisons expected in current TKDE and KDD research
\citep{wu2024supporting,zhang2024oagbench}.  These methods address broader problems
than the present metadata intervention, but they motivate the inclusion of a supervised
LambdaRank fusion rather than comparison only with fixed linear weights.

\subsection{Bibliographic signals and bias}

Citation count is an attention measure, not a direct measure of scientific quality.
It varies with field, publication age, language, and network visibility; responsible
metric use requires multiple indicators and explicit context
\citep{hicks2015leiden}.  A current citation snapshot is especially problematic when
the target itself is a citation edge.  Our confirmatory target is instead human graded
relatedness.  Even there, current metadata remains observational and time-dependent,
so the experiment supports neither a causal claim nor a historical simulation.

\subsection{Selective decisions and risk}

Selective prediction evaluates a model jointly with a reject option and studies risk
as coverage changes \citep{geifman2017selective,geifman2019selectivenet}.  High-
confidence policy improvement similarly compares a proposed action with a baseline
under explicit assumptions \citep{thomas2015hcpi}.  BiblioGuard borrows the operational
idea of a pessimistic score but does not inherit a theorem: similarity-selected
neighbours are dependent and the weighted standard error is heuristic.  We therefore
calibrate a threshold on a separate split and call it an empirical operating rule.

\section{Problem formulation}

For seed paper \(q\), let \(D_q\) be its fixed candidate pool, \(s_0(q,d)\) the
SPECTER2 cosine score, and \(y_{qd}\in\{0,1,2\}\) the locked relevance grade.  The
content ranking is \(\pi_0(q)\).  An action \(a\in\mathcal{A}\) changes scores but not
the candidate pool.  Its query-level effect is
\begin{equation}
  \Delta_q(a)=\operatorname{NDCG@10}(\pi_a(q),y_q)
  -\operatorname{NDCG@10}(\pi_0(q),y_q).
\end{equation}
A policy chooses \(a\) or abstains.  Abstention has effect zero.  The primary estimand
is the mean policy effect over every locked query, not merely active queries.

Three validity constraints follow.  First, action decisions cannot use locked labels.
Second, action prediction must remain fixed while a gate threshold is swept; otherwise
the curve confounds action and coverage.  Third, methods must be compared at the same
coverage as well as at their native calibrated operating points.

\section{Method}

\subsection{Pre-specified metadata actions}

Within each candidate pool, SPECTER2 score, \(\log(1+\text{citations})\), and year are
min--max scaled.  A missing value is replaced by the observed pool median before
scaling; a fully missing or constant field maps to zero.  The action family contains
citation, recency, and equally balanced interpolations at weights 0.15 and 0.30, plus
reciprocal-rank fusion of SPECTER2 with citation or recency rank at \(k=60\).  No
action or weight is added after test evaluation.  Equal metadata values receive
midranks in reciprocal-rank fusion, so a constant or fully missing metadata field
cannot create an identifier-driven ordering.

\subsection{Local action prediction}

Word (1--2 gram) and character (3--5 gram) TF--IDF features are fitted on training
seed titles.  A target uses \(k=\lfloor\sqrt{n_{\mathrm{train}}}\rfloor\) nearest
training seeds.  Negative cosine similarities are clipped and the remaining
similarities are normalised to weights \(w_i\).  For action \(a\),
\begin{align}
  \widehat\mu_a &= \sum_i w_i\Delta_i(a),\\
  \widehat v_a &= \frac{\sum_i w_i(\Delta_i(a)-\widehat\mu_a)^2}
    {1-\sum_i w_i^2},\\
  \widehat{\operatorname{SE}}_a &=
    \sqrt{\widehat v_a\sum_iw_i^2}.
\end{align}
The action predictor is \(\arg\max_a\widehat\mu_a\).  Importantly, it is fixed before
gating.  The confidence score is
\(g=\widehat\mu_{a^*}-1.645\widehat{\operatorname{SE}}_{a^*}\).

\subsection{Calibration and comparators}

On the independent calibration split, the threshold is the least restrictive value
whose one-sided 95\% Clopper--Pearson upper bound on conditional harm is at most
0.20 with at least 30 active queries \citep{clopper1934}.  Among eligible thresholds,
coverage is maximised.  If none is eligible, the policy abstains everywhere.

Comparator selectors are (i) local kNN mean without the uncertainty penalty,
(ii) a multi-output histogram gradient-boosted gain predictor, and (iii) a regressor
that gates the single globally best training action.  A LightGBM LambdaRank model
\citep{ke2017lightgbm} is a separate ranking baseline trained on BM25, BGE, SciNCL,
SPECTER2, log-citation, and year features.  It is fitted for 500 trees on training
labels only, with linear graded gains $[0,1,2]$; calibration and locked labels are
not used for fitting or early stopping.

\section{Locked experimental protocol}

\subsection{Data and split}

The pinned SciRepEval RELISH input contains 3,190 seed papers,
\CandidatePairs{} judged pairs, and
\UniqueDocuments{} unique documents.  Exact normalised seed titles form groups.
The SHA-256 of \code{relish-v1|group} assigns buckets 0--1 to training, bucket 2 to
calibration, and buckets 3--9 to the locked test.  This produced \TrainQueries{},
\CalibrationQueries{}, and \TestQueries{} seed papers, with no identical normalised
title crossing splits.

A pre-test character 3--5-gram TF--IDF audit found 17 cross-split title pairs with
cosine similarity at least 0.70 and four locked queries with similarity at least 0.80
to a training or calibration title.  Membership was not changed after this audit.  We
instead froze those four identifiers and report a secondary analysis excluding them,
while retaining all 2,240 queries in the primary result.

The test labels have three levels: irrelevant, somewhat relevant, and relevant.  We
retain every query.  Candidate order is not used as a feature, and score ties are
broken by the dataset document identifier.  The raw input and label Parquet files are
pinned to immutable Hugging Face revisions and checksummed.

\subsection{Metadata and model provenance}

Semantic Scholar metadata are requested only by exact, non-missing CorpusId.  Raw
responses, request fields, retrieval times, missing records, title checks, and the
snapshot hash are stored.  The 1,788 unique documents whose CorpusId is encoded as
missing are not string-matched to external records.  This avoids the silent entity
resolution errors present in the earlier repository.

All neural models are resolved to immutable Hugging Face commit hashes.  Documents
are encoded locally from title and abstract on an NVIDIA RTX 4060; SciNCL uses its
required tokenizer separator between the two fields, BGE uses a space, and SPECTER2
uses its base tokenizer separator.  No old embedding
cache is accepted.  Package versions, CUDA build, hardware, input hashes, score
hashes, policy decisions, and Git commit are included in the run manifest.

\subsection{Outcomes and inference}

The primary metric is NDCG@10 with linear graded gain, matching the
trec\_eval/pytrec\_eval convention used by SciRepEval \citep{jarvelin2002}.  Secondary
metrics are NDCG@20, full-list NDCG, MAP@10, Recall@50, and Precision@10.  MAP@10
uses the total number of relevant candidates in its denominator, as in
\code{map\_cut.10}; it is not renormalised when more than ten candidates are relevant.
At the calibrated operating point and coverage
budgets 10\%, 25\%, 50\%, 75\%, and 100\%, we report overall mean gain, active-query
mean gain, \(P(\Delta<0\mid\mathrm{active})\),
\(P(\Delta\le -0.05\mid\mathrm{active})\), and mean negative shortfall
\(E[\max(0,-\Delta)\mid\mathrm{active}]\).  Selective risk is the latter conditional
expectation as queries are included in descending confidence order; AURC is its
trapezoidal area.

The mean primary effect receives a 10,000-replicate paired query bootstrap
interval.  A two-sided paired randomisation test uses 100,000 sign flips.  Secondary
method comparisons use Holm adjustment \citep{holm1979}.  These analyses quantify
sampling uncertainty over the fixed fitted policy; they do not establish a
distribution-free safety guarantee.

\subsection{Retrospective audit}

SCIDOCS, SciFact, NFCorpus, and TREC-COVID are not reused as confirmatory tests.  The
repository audit found that SCIDOCS current citation counts were coupled to its
citation-prediction target, official tests had influenced method choices, SciFact and
NFCorpus contained cross-split duplicate query texts, and cache files lacked complete
query-order and metadata hashes.  The earlier ``fuzzy gate'' claims are removed
because neither an implementation nor result lineage exists on any repository branch.

\section{Results}

\subsection{Retriever and learned-fusion baselines}

\IfFileExists{generated/table_retrievers.tex}{
  \input{generated/table_retrievers.tex}
}{
  \begin{center}\fbox{Retriever table is generated only from the completed locked run.}\end{center}
}

The comparison uses the same candidate documents and labels for every system.  We
interpret the strongest locked baseline descriptively; it does not change the
pre-specified SPECTER2 backbone or the primary comparison.

\subsection{Calibrated operating points and matched coverage}

\IfFileExists{generated/table_policy.tex}{
  \input{generated/table_policy.tex}
}{
  \begin{center}\fbox{Policy table is generated only from the completed locked run.}\end{center}
}

\IfFileExists{generated/figure_risk_coverage.pdf}{
  \begin{figure}[t]
    \centering
    \includegraphics[width=0.82\linewidth]{generated/figure_risk_coverage.pdf}
    \caption{Conditional negative-shortfall risk as coverage increases.  Each selector
    keeps its action predictor fixed; only its confidence threshold changes.}
    \label{fig:riskcoverage}
  \end{figure}
}{}

\subsection{Primary inference and failure analysis}

At its frozen operating point, BiblioGuard changes mean NDCG@10 by \PrimaryGain{}
with 95\% interval \PrimaryCI{} and paired-randomisation \(p=\PrimaryPValue\).
Coverage is \PrimaryCoverage{}.  The claim made from this result follows the stopping
rule stated before label access: superiority is asserted only if the policy exceeds the
content baseline and remains competitive with the fixed-action and learned selectors
under equal coverage.

\IfFileExists{generated/table_actions.tex}{
  \input{generated/table_actions.tex}
}{}

\section{Discussion}

\subsection{What the locked result does and does not show}

The experiment isolates current bibliographic metadata on a human relatedness target.
It can show whether such signals improve a fixed 2026 recommendation snapshot under
the stated candidate pool and development distribution.  It cannot show that citation
count caused relevance, that the same policy would have worked historically, or that
offline NDCG implies downstream researcher benefit.

The strongest simple comparator is central to interpretation.  If a global action has
the best mean effect, the main finding is that a popularity/recency prior is useful on
this collection, not that local routing is necessary.  If LambdaRank is best, the
selective layer may still reduce exposure but is not the best ranker.  If calibration
forces complete abstention, it prevents an unsupported intervention but adds no
retrieval value.  These are substantive outcomes rather than presentation problems.

\subsection{Why the earlier positive result is withdrawn}

The original SCIDOCS number mixed a post-benchmark citation snapshot with a citation-
prediction outcome.  The same test collection was then used for feature, neighbour,
and penalty ablations.  Its reported score is numerically reproducible from stored
arrays, but it is not clean confirmatory evidence.  Reclassifying it as retrospective
is necessary even though it weakens the paper's headline.

\subsection{Responsible use}

Citation-based reranking can amplify established venues, older papers, dominant
languages, and already visible groups.  The system should expose which signal changed
a ranking and permit users to disable metadata.  Missing metadata must remain missing;
title-based substitution without verified identifiers is inappropriate for a decision
study.  Abstention limits the number of interventions but does not remove bias among
the interventions that remain.

\section{Limitations}

First, the locked confirmatory collection is biomedical; other domains in the old
study are retrospective.  Second, RELISH uses candidate pools assembled by earlier
retrieval systems, so results do not evaluate full-corpus recall.  Third, current
citation counts and years are an explicitly dated snapshot.  Fourth, the local
standard error does not model dependence among related training papers.  Fifth, the
calibration rule controls only an empirical split-specific upper bound and may choose
zero coverage.  Sixth, 2,240 test seeds quantify offline ranking, not user satisfaction
or discovery diversity.  Seventh, four locked titles have high character similarity
to development titles; the frozen sensitivity analysis addresses but cannot eliminate
all topic dependence.  Eighth, model and metadata providers may later change, even
though the resolved artefacts and hashes used here are frozen.

\section{Reproducibility}

The public repository contains the locked protocol, exact data and model revisions,
fully pinned environment, phased command-line pipeline, leakage tests, frozen
decision hash, per-query locked outcomes, generated tables, and manuscript source.
Aligned candidate scores, citation/year vectors, and the compressed metadata snapshot
are stored with the release; large public inputs and document embeddings are regenerated.
A
verification command checks every published number against the machine-readable
result file.  A v3.1 erratum records metric, preprocessing, and integrity corrections
found during code audit; it was committed before locked-test access.  The original
protocol commit predates labelled RELISH evaluation, and the freeze
manifest explicitly records \code{test\_labels\_consumed: false}.

\section{Conclusion}

We replaced a favourable but confounded four-benchmark narrative with a locked,
leakage-aware test of bibliographic intervention.  The revised contribution is not
that citation and recency are universally useful, nor that a heuristic lower score is
a safety theorem.  It is a falsifiable selective-decision protocol, a fair comparison
against modern retrieval and learned fusion, and an auditable account of query-level
benefit and harm.  The empirical conclusion is determined by the frozen RELISH run
rather than selected after it.

\section*{Declarations}

\textbf{Funding.} The author received no specific funding for this work.

\textbf{Competing interests.} The author declares no competing interests.

\textbf{Ethics.} The computational study uses public, de-identified benchmark data.
RELISH's original collection and consent procedures are described by
\citet{brown2019relish}.

\textbf{Data and code availability.} Code, protocols, hashes, and non-restricted
results are available in the public PaperPilot-Reproduction repository.  The final
commit identifier is inserted after the verified release; no archival DOI is claimed
until one exists.

\textbf{Author contributions.} Zichen Feng: conceptualisation, methodology, software,
validation, formal analysis, investigation, writing---original draft, and
writing---review and editing.

\bibliographystyle{plainnat}
\bibliography{references}
\end{document}
